\documentclass[11pt]{article}
\usepackage[margin=1in]{geometry}
\usepackage{amsmath,amssymb,amsthm}
\usepackage[hidelinks]{hyperref}

\newtheorem*{theorem}{Literature Answer}

\title{Negative Literature Answer to Question 5.3 of arXiv:2105.09203}
\author{FA/Banach run packet}
\date{2026-06-18}

\begin{document}
\maketitle

\section*{Status}

This is a literature-already-answered packet. It records that Question 5.3
from Ansorena--Bello--Wojtaszczyk \cite{ABW} was answered negatively in the
later paper of Albiac--Ansorena--Berasategui \cite{AAB}. It is not an
original proof from this run.

\section*{Original Problem}

In the open-problems section of \cite{ABW}, the third question, cited later as
Question 5.3, asks:

\begin{quote}
Let \(\mathcal B\) be a bidemocratic basis of \(\ell_2\). Are there
\(1<q<r<\infty\) such that
\[
\ell_{2,q}\stackrel{\mathcal B}{\hookrightarrow}\ell_2
\stackrel{\mathcal B}{\hookrightarrow}\ell_{2,r}?
\]
\end{quote}

\section*{Supporting Answer}

The supporting paper \cite{AAB} constructs bidemocratic bases with large
unconditionality parameters. In Corollary 4.6 it obtains, for separable Hilbert
space in particular, a bidemocratic total basis \(\mathcal B\) with
\[
\tilde k_m[\mathcal B]\approx \log(1+m).
\]
The following remark observes that any basis squeezed between suitable Lorentz
spaces \(d_q(w)\) and \(d_r(w)\), with \(q>1\) or \(r<\infty\), must satisfy a
strictly smaller logarithmic growth estimate for \(k_m\). It then explicitly
states that, for Hilbert space, the constructed basis is a counterexample that
solves \cite[Question 5.3]{ABW} negatively.

\begin{theorem}
Question 5.3 of \cite{ABW} has a negative answer. There exists a bidemocratic
basis of \(\ell_2\) for which no exponents \(1<q<r<\infty\) give the
embeddings
\(\ell_{2,q}\stackrel{\mathcal B}{\hookrightarrow}\ell_2
\stackrel{\mathcal B}{\hookrightarrow}\ell_{2,r}\).
\end{theorem}

\begin{proof}[Literature identification]
The basis supplied by Corollary 4.6 of \cite{AAB}, with the ambient space
chosen to be Hilbert space, is bidemocratic and has logarithmic
unconditionality growth. The subsequent remark in \cite{AAB} says that a basis
with the Lorentz squeezing requested in Question 5.3 of \cite{ABW} would have
to satisfy
\[
k_m[\mathcal B]\lesssim (1+\log m)^{1/q-1/r}.
\]
For \(1<q<r<\infty\), the exponent \(1/q-1/r\) is strictly less than \(1\),
contradicting the logarithmic lower growth of the constructed basis. The
supporting authors record exactly this implication and state that the Hilbert
space instance is a negative solution of \cite[Question 5.3]{ABW}.
\end{proof}

\section*{Scope Limitations}

This packet answers only Question 5.3 of \cite{ABW}. The preceding questions
in that open-problems section, concerning quasi-greedy bases in
superreflexive spaces and dual quasi-greediness in \(\ell_p\), are not claimed
here.

\section*{Verification Notes}

The source and supporting papers are distinct. The supporting paper explicitly
cites \cite[Question 5.3]{ABW}, so this belongs in
\texttt{literature\_already\_answered}. The local duplicate search covered
arXiv ids \texttt{2105.09203} and \texttt{2205.09478}, plus phrases
\texttt{Question 5.3}, \texttt{bidemocratic basis}, \texttt{ell\_2},
\texttt{Lorentz}, and \texttt{sparse approximation using new greedy-like
bases}.

\begin{thebibliography}{9}

\bibitem{ABW}
J. L. Ansorena, G. Bello, and P. Wojtaszczyk,
\emph{Lorentz spaces and embeddings induced by almost greedy bases in
superreflexive Banach spaces},
arXiv:2105.09203.

\bibitem{AAB}
F. Albiac, J. L. Ansorena, and M. Berasategui,
\emph{Sparse approximation using new greedy-like bases in superreflexive
spaces},
arXiv:2205.09478.

\end{thebibliography}

\end{document}
